\section{The affine-packet programme: source reconstruction and exact repairs}
\label{sec:affine-packet-source-audit}

This chapter reconstructs the affine-packet mathematics in the 3,021-line modern
Chatnotes export \emph{Implications of Affine Formula} and of its immediate
61-page local antecedent, the OpenAI technical note \emph{Affine Packets,
Cyclic Orbit Packets, Dyadic Dilation, Lambert--Mahler Series, Hydra/Numen
Structure, and Exact Rational Periodic Orbit Enumeration in the Odd-Step
Collatz Algebra}.  Neither local item is used as a substitute for published
literature.  The purpose of the reconstruction is to identify exactly which
coordinates are classical, which maps have now been proved, which finite
calculations reproduce, and which statements in the antecedent require
repair.

The raw export has SHA-256
\begin{center}
\texttt{abb1abbb948dc140394de3d35f0efbde9b3facdf718c24b065b8d0adb951366e}
\end{center}
and 3,021 physical lines.  The local technical note has SHA-256
\begin{center}
\texttt{13f49961a60b870fa7a6f0f646758b4d72c1a4e4267f347606cc8a3238d552f1}.
\end{center}
All page and line locators below refer to those exact manifestations.

\subsection{The local sources}

The raw export has the following chronological organization.  Physical line
numbers refer to the export itself, not its generated attachment labels.
\begin{center}
\begin{tabular}{p{0.16\textwidth}p{0.72\textwidth}}
\toprule
Physical lines & source content\\
\midrule
1--613 & First synthesis of the attached local technical note; source names
and theorem numbers here are routing prompts, not independent evidence.\\
614--1373 & Recheck of the fixed-return algebra, cyclic-shift bridge,
primitive necklaces, signed rational counts, packet envelopes, and the
length-four argument.\\
1375--2001 & Length-five argument; cycle product identity; exact
minimum-point strip; reduction of period eight to $(m,A)=(8,13)$.\\
2003--2679 & Further derivation of the strip and candidate packet list.\\
2680--3020 & First-letter residue recursion, the claimed
$R_{8,13}(233)$ calculation, the 99 cyclic classes, and exhaustive candidate
packet scans through $m=16$.\\
3021 & Exporter footer; no mathematical content.\\
\bottomrule
\end{tabular}
\end{center}

The export refers to attached CSV, Markdown, and summary tables, but those
objects are not embedded in the 3,021-line file.  Accordingly no calculation
was admitted from a link label.  The 99-class table was independently
recreated as
\path{research_companion/data/period_8_A_13_classes.csv}, and every row is
checked by the portable certificate.

The antecedent develops affine blocks, translations, cyclic transport,
necklace counting, and several analytic constructions.  Its Theorems~10.1
and 10.3 supply the cyclic-transport and primitive-necklace arguments examined
below.  The indexed source TeX and a second, 62-page PDF agree with the
controlling PDF in the twelve statement/proof blocks collated for this
chapter; the observed differences there are populated cross-references and
repagination.  This is a passage-level correspondence, not an identification
of every unexamined statement or of the exact producing builds.

\subsection{Published antecedents and the exact coordinate bridge}

The core return formula predates the local corpus.

\begin{itemize}[leftmargin=2em]
\item B\"ohm--Sontacchi, printed pp.~263--264, give the formal affine
composition and fixed-return candidate for a prescribed parity pattern
\cite{BohmSontacchi1978}.

\item Crandall, printed pp.~1287--1290, gives the accelerated affine cycle
formula, the minimum-orbit estimate, and a cycle-length exclusion based on the
then-current verification range \cite{Crandall1978}.  Thus the product/minimum
strip in raw lines 1581--1620 and 1829--1908 is a direct form of a classical
cycle inequality, not a new structural principle.

\item Lagarias, printed pp.~34--39, proves that every binary parity word gives
a unique rational periodic point, displays its numerator and denominator,
defines irreducibility by cyclic rotation, and counts irreducible binary
necklaces \cite{Lagarias1990}.  In particular, a new vocabulary for the same
numerator cannot enlarge the theorem.

\item Urata, printed pp.~8--11, writes the odd-step exponent numerator,
proves the periodic-exponent/periodic-orbit correspondence, and asks exactly
when
\[
 \frac{C(p)}{2^{a_1+\cdots+a_m}-3^m}
\]
is an integer \cite{Urata2003}.

\item Laarhoven--de Weger, source lines 303--321, relate rational cycles to
binary necklaces/Lyndon words and to primitive $3n+b$ cycles
\cite{LaarhovenDeWeger2012}.  Their statements retain the source defects and
scope boundaries recorded in Chapter~\ref{sec:literature-spine}; they are not
used to erase the primitive-denominator conditions in Lagarias.
\end{itemize}

For a positive exponent packet $p=(a_1,\ldots,a_m)$, $m\ge1$, of
total $A\ge m$, define
\[
 E(p)=1,0^{a_1-1}\cdots1,0^{a_m-1}.
\]
This word uses Lagarias's shortened map $T$, whose odd branch is
$(3x+1)/2$ and whose even branch is $x/2$, on the rationals with odd
denominator.  It is a bijection from compositions of $A$ into $m$ positive parts to
binary words of length $A$, weight $m$, and initial bit one.  If
$S_j=a_1+\cdots+a_j$, the ones occur at positions
$S_0,\ldots,S_{m-1}$, so Lagarias's numerator is exactly
\[
 N(E(p))=\sum_{j=0}^{m-1}3^{m-1-j}2^{S_j}=C(p).
\]
Packet rotation by $j$ letters is not one-bit rotation: it is binary rotation
by the prefix length $S_j$.  The rotation comparison is the groupoid isomorphism
\[
 (p,j)\longmapsto(E(p),S_j(p)),
 \qquad E(\sigma^jp)=\rho^{S_j(p)}E(p),
\]
with inverse given by cyclic gaps between successive ones.  The composition
identity
\[
 S_{j+k\bmod m}(p)
 \equiv S_j(p)+S_k(\sigma^jp)\pmod A
\]
proves functoriality.  It also proves that primitive exponent packets and
primitive binary parity words are the same periodic objects in different
coordinates.  The complete domain, codomain, inverse, fibres, and stabilizers
are proved in the companion's theorems \emph{Packet--parity coordinate
isomorphism} and \emph{Exact cyclic groupoid morphism}, in its section
\emph{Affine packets, parity necklaces, and exact residue automata}.
Here the acting groups are $\mathbb Z/m\mathbb Z$ and
$\mathbb Z/A\mathbb Z$; rotation arrows with equal targets remain distinct.
For $p=u^r$, with $u$ primitive of length $d$ and sum $B$, the stabilizer
map is $kd\bmod m\mapsto kB\bmod A$ for $0\le k<r$.

For the full map $F$, with odd branch $3x+1$ and even branch $x/2$, the
corresponding word is instead
\[
 E_F(p)=1,0^{a_1}\cdots1,0^{a_m}.
\]
It has length $A+m$.  Inserting a zero after each one of $E(p)$, or deleting
that zero from $E_F(p)$, gives mutually inverse pointed-word maps.
The exact odd-return times are
\[
 T^{S_j}(x_p)=\Phi^j(x_p)=F^{S_j+j}(x_p).
\]
For the primitive root $u$, the exact periods are $d$ odd returns, $B$
shortened steps, and $B+d$ full steps.  The companion's theorem
\emph{Exact comparison of three clocks} proves the intermediate values
and minimal periods, rather than identifying the clocks by terminology.

\subsection{Cyclic affine transport and exact valuations}

The antecedent already proves the affine arrow explicitly in
Theorem~10.1, physical p.~36: $g_{a_1}(x_p)$ is fixed by $g_{\sigma p}$,
whose fixed point is unique.  Raw lines 959--1003 give the corresponding
numerator calculation.  The additional check needed to identify the actual
odd-return branch is its exact valuation.  With
\[
 D=2^A-3^m,
 \qquad x_p=\frac{C(p)}D,
 \qquad \sigma p=(a_2,\ldots,a_m,a_1),
\]
the first-letter recurrences, with $C(\varnothing)=0$ for the $m=1$
tail, give
\[
 3C(p)+D=2^{a_1}C(\sigma p),
 \qquad 3x_p+1=2^{a_1}x_{\sigma p}.
\]
Because $C(\sigma p)$ and $D$ are odd,
\[
 \nu_2(3x_p+1)=a_1,
 \qquad \Phi(x_p)=x_{\sigma p}.
\]
This proves an actual odd-step morphism; it is not merely equality of two
formal affine expressions.  Since $\gcd(D,6)=1$, both $2$ and $3$ are
units modulo $D$.  The first identity therefore gives
$D\mid C(p)$ if and only if $D\mid C(\sigma p)$, proving invariance
under cyclic rotation.

The exact positive-cycle criterion is therefore:
\[
 \left.
 \begin{array}{c}
 p\text{ primitive},\\
 2^A>3^m,\\
 2^A-3^m\mid C(p)
 \end{array}
 \right\}
 \quad\Longleftrightarrow\quad
 \text{a positive odd integral orbit of exact odd-step period }m
\]
with exponent necklace $[p]$.  If $p$ is imprimitive, the same rational point
belongs to the orbit of its primitive root; a length-$m$ return is not thereby
an exact period-$m$ orbit.

\subsection{Exact repairs to the 61-page local antecedent}

The following are mathematical repairs, not stylistic preferences.

\begin{sourceissue}[Periodic-tail calculation, physical p.~30]
\label{issue:affine-periodic-tail-expansion}
Theorem~9.2 asserts that an integral orbit with an eventually repeated
exponent block $p$ is already fixed under the block map from the beginning
of that exponent tail.  This assertion is correct.  Write
$D=2^A-3^m$, $C=C(p)$, and $z_j=Dy_j-C$ for the integral block samples
$y_{j+1}=(3^my_j+C)/2^A$.  The required calculation is
\[
 z_{j+1}
 =\frac{3^mDy_j+DC-2^AC}{2^A}
 =\frac{3^m(Dy_j-C)}{2^A}
 =\frac{3^m}{2^A}z_j.
\]
The printed intermediate numerator replaces $DC$ by $3^mC$; the
following line already gives the correct result.  For $p=(2)$ and
$y_j=1$, the incorrect line equals $1/2$, whereas $z_{j+1}=0$.

After the displayed repair, $2^{jA}z_j=3^{jm}z_0$ and
$\gcd(2^{jA},3^{jm})=1$ give $2^{jA}\mid z_0$ for every $j$.
Since $A\ge1$, a nonzero integer $z_0$ cannot have this divisibility:
choose $j$ with $2^{jA}>|z_0|$.  Hence $z_0=0$ and $y_j=C/D$ for
all $j$.  The intermediate odd states repeat because the map is
deterministic.  This proves exact repetition, not only convergence.
The hypothesis that the exponent tail repeats is retained; no eventual
periodicity of an arbitrary orbit is inferred.

The same intermediate error occurs in the indexed 62-page PDF, p.~31,
and in the indexed TeX at line~1977.  These particular passages were
collated directly; no whole-version identity is needed.
\end{sourceissue}

\begin{sourceissue}[Complex Abel coordinate, physical pp.~11--12]
On a simply connected invariant domain, arbitrary choices of a holomorphic
logarithm need not satisfy
$\operatorname{Log}(\chi w)=\log\chi+\operatorname{Log}(w)$ literally; the
two sides can differ by $2\pi i k$.  The transport statement requires a
branch-compatible choice, or an explicit domain argument forcing $k=0$.
Later portions of the note impose compatibility, but Theorem~4.2 as printed
does not.
\end{sourceissue}

\begin{sourceissue}[Concentration calculation, physical p.~18]
For $f(t)=t^2\log^22$ and $x\ge y\ge2$,
\[
 f(x+1)+f(y-1)-f(x)-f(y)
 =2\log^22\,(x-y+1)>0.
\]
The stated concentration conclusion is correct, but the displayed
application in Corollary~6.5 reverses the sign in its rearrangement.  The
identity above is the repaired proof.
\end{sourceissue}

\begin{sourceissue}[Maximizer multiplicity, physical p.~19]
At the boundary $A=m$ there is only the all-ones packet, so the maximizing
word has multiplicity one.  More generally, the multiplicity is one when
$m=1$ or $A=m$, and is $m$ when $m\ge2$ and $A>m$, when the unique large
entry can occupy any of the $m$ positions.
Corollary~6.7 omits this boundary case.
\end{sourceissue}

\begin{sourceissue}[The $m=1$ proof boundary, physical pp.~21--23]
\label{issue:affine-singleton-average}
The displayed formula in Theorem~6.10 remains correct for $m=1$, but the
printed proof sets $r=m-2$, whereas its auxiliary lemma on p.~20 assumes
$r\ge0$.  For $m=1$ there is exactly one packet, $(A)$, and hence
\[
 \mathbb E[a_1^2]=A^2=\frac{A(2A-1+1)}{1(1+1)},\qquad
 \mathbb E[Q]=(A\log2-\log3)^2.
\]
These are precisely the displayed formulas at $m=1$, completing that case.
\end{sourceissue}

\begin{sourceissue}[Pointed periodic words, physical p.~37]
\label{issue:affine-pointed-words}
The last sentence of the proof of Theorem~10.3 incorrectly makes equality
of infinite pointed periodic words equivalent to cyclic conjugacy of their
primitive roots.  For example, $(1,2)^\infty\ne(2,1)^\infty$.
The correct statement is equality of their primitive roots as pointed words.
Indeed, extend a common periodic sequence to all integer indices.  Its
periods form a subgroup of $\mathbb Z$, so their least positive element
divides every period; each word is a power of that common initial block.
Primitive words therefore equal that block.  In the source proof the base
points have already been aligned, so this stronger conclusion supplies the
needed implication.  Passing back to unpointed orbits gives cyclic conjugacy,
as stated in the theorem.
\end{sourceissue}

\begin{sourceissue}[One-letter Hamiltonian domain, physical p.~55]
\label{issue:affine-H1-domain}
The source Hilbert space is $\ell^2(\mathbb N_{\ge1})$, and the multiplier
is $u(a)=a\log2-\log3$.  The self-adjoint operator is
\[
 \begin{split}
 \mathcal D(H_1)&=\left\{\xi\in\ell^2(\mathbb N_{\ge1}):
       \sum_{a\ge1}|u(a)|^2|\xi_a|^2<\infty\right\},\\
 (H_1\xi)_a&=u(a)\xi_a.
 \end{split}
\]
To verify self-adjointness, test the adjoint identity on each basis vector:
an adjoint-domain vector $\eta$ must have adjoint image
$(u(a)\eta_a)_a\in\ell^2$, and hence belongs to the displayed domain.
Conversely this condition suffices by the Cauchy--Schwarz inequality.
Thus the adjoint has exactly the same domain and action.  Truncating at
$a\le N$ converges both in $\ell^2$ and after application of $H_1$;
finite-support vectors form a core, giving an essentially self-adjoint
restriction, not a second unspecified self-adjoint operator.

The resolvent at $i$ is the diagonal operator with entries
$(u(a)-i)^{-1}$.  Since $|u(a)/(u(a)-i)|\le1$, this inverse maps
$\ell^2$ into the stated operator domain.  Its entries tend to zero, and
its truncations converge in operator
norm, so the resolvent is compact.  Finally, for every $\beta>0$,
\[
 \operatorname{Tr}(e^{-\beta H_1})
 =\sum_{a\ge1}3^\beta2^{-\beta a}
 =\frac{3^\beta}{2^\beta-1},
\]
as in Theorem~19.1.  Specifying the domain completes this operator statement
without changing its spectrum or partition function.
\end{sourceissue}

The note also contains unresolved bibliography/cross-reference markers and
named analogies that are not accompanied by typed maps.  Those analogies are
not mathematical identifications and are not propagated.  No theorem is
discarded merely because it uses local vocabulary; each formula is instead
crosswalked through its explicit coordinates.

\subsection{What reproduces from the raw Chatnotes}

\subsubsection{Signed fixed-content necklace counts}

Raw lines~1041--1083 give a valid fixed-label refinement of the classical
primitive-necklace count.  For $m\geq1$, $A\geq m$, let $W(m,A)$ count the
primitive words in $\Pi_{m,A}$.  Unique primitive-root decomposition and
M\"obius inversion give
\[
 W(m,A)=\sum_{d\mid\gcd(m,A)}\mu(d)
          \binom{A/d-1}{m/d-1}.
\]
The cyclic action is free on primitive words, so the number of rational
$\Phi$-orbits of exact odd-return period $m$ at this packet label is
\[
 N_{\mathrm{orb}}(m,A)
 =\frac{W(m,A)}m
 =\frac1A\sum_{d\mid\gcd(m,A)}\mu(d)
          \binom{A/d}{m/d}.
\]
The second equality is the fixed-Hamming-weight binary-necklace formula.
Stanley's multivariate Lyndon series enumerates primitive necklaces by their
colour multiplicities; coefficient extraction at binary content $(m,A-m)$
gives exactly this expression
\cite[Exercise~7.89(a), equations~(7.149)--(7.150), and the solution on
printed p.~555]{Stanley2024}.  Lagarias gives the unrestricted primitive
binary-necklace and rational-cycle coordinates, and Laarhoven--de Weger
restate the unrestricted count in Lyndon-word language
\cite[pp.~38--39]{Lagarias1990}\cite{LaarhovenDeWeger2012}.  The new step in
the present reconstruction is not the enumeration formula: it is the exact
transport through the packet--parity groupoid followed by the
Collatz-denominator sign partition.

Every point belonging to the label $(m,A)$ is
$C(\sigma^jp)/(2^A-3^m)$ with positive numerator.  Hence the whole orbit is
positive when $2^A>3^m$ and negative when $2^A<3^m$; equality is impossible.
Multiplication by the relevant sign indicator gives the two signed orbit
counts, and multiplication by $m$ gives the two point counts.  The companion
states and proves all three typed counts separately.  Raw line~1083 instead
calls every word a periodic point.  The exact statement is that a primitive
word determines a distinguished point, while its $m$ rotations determine the
$m$ distinct points of one exact-period orbit.  When $\gcd(m,A)=1$, every
word is primitive and the orbit and point counts are respectively
$m^{-1}\binom{A-1}{m-1}$ and $\binom{A-1}{m-1}$.

\subsubsection{The analytic packet family}

The raw series at lines~432--475 are substantive content, not merely
motivation.  In the source coordinates $c(p)=C(p)/2^A$, put
\[
 S_m(A)=\sum_{p\in\Pi_{m,A}}c(p),\qquad
 X_m(A)=\sum_{p\in\Pi_{m,A}}\frac{C(p)}{2^A-3^m},\qquad
 F_m(q)=\sum_{A\ge m}S_m(A)q^A.
\]
The source's Theorems~16.2--16.5 give, with
$u=q/(1-q)$ and $v=q/(2-q)$,
\[
 F_{m+1}=v(3F_m+u^m),\quad F_1=v,\qquad
 F_m=\sum_{k=0}^{m-1}3^{m-1-k}u^kv^{m-k}.
\]
These identities follow by appending a letter, using
$c(p,a)=2^{-a}(3c(p)+1)$, and summing a finite convolution at every
coefficient.  They supply the stated two-state linear representation.
The companion section \emph{Packet fixed-point series and their dyadic
poles} reconstructs this family with its full proofs.

\begin{sourceissue}[Generalized dilation equation, physical p.~51]
\label{issue:affine-mahler-domain}
Theorem~17.1 proves the coefficientwise identity
\[
 X_m(2q)-3^mX_m(q)=F_m(2q),\qquad
 X_m(q)=\sum_{A\ge m}X_m(A)q^A,
\]
and its uniqueness, because $2^A-3^m\ne0$.  The additional claim that
the defining series therefore give this identity throughout $|q|<1$
fails when $m\ge2$.  Already for $m=2$,
\[
 X_2(A)=\frac{2^A+3A-5}{2^A-9}\longrightarrow1,
\]
so $X_2(2q)$ diverges at $q=3/4$.  The phrase also occurs in the
indexed PDF, p.~52, and in the indexed TeX at line~3376; it is not
present in the raw Chatnotes statement.

The exact common Taylor-series domain is $|q|<1/2$ for $m\ge2$;
for $m=1$ it is $|q|<1$.  Indeed, for large $A$ the total $X_m(A)$
is positive and bounded above by a constant times
$\binom{A-1}{m-1}$.  If $m\ge2$, a word ending in $1$ contributes
at least $1/2$, so the radius of $X_m$ is exactly one.  For $m=1$,
$X_1(A)=(2^A-3)^{-1}$ has radius two.  The companion gives the
coefficient bounds and the inverse of the formal dilation map explicitly.
\end{sourceissue}

The source already continues $X_1$ meromorphically with simple poles at
$2,4,8,\ldots$.  The companion extends this calculation to the whole
packet family: for $m\ge2$, $X_m$ has a pole of order $m-1$ at $1$
and a pole of order $m$ at every $2^j$, $j\ge1$, with no others.
The construction subtracts a finite coefficient prefix and sums scaled
copies of the remaining rational function with a uniform geometric tail
bound.  It proves a meromorphic identity, not convergence of the original
series outside its radius.  No claim of priority over the broader
literature accompanies this extension of the local source calculation.

\subsubsection{The finite cycle calculations}

The length-four and length-five arguments at raw lines 1209--1365 and
1567--1825 are arithmetically correct after the classical lineage above is
made explicit.  For length five the strip forces $A=8$; the denominator is
$13$; and the seven cyclic representatives have numerator residues
\[
 3,6,12,4,10,6,7\pmod{13}.
\]
None is integral.

The packet residue recurrence at raw lines 2683--2790 is also exact.  For a
positive odd modulus $q$,
\[
 R_{m,A}(q)=\{C(p)\bmod q:p\in\Pi_{m,A}\}
\]
satisfies
\[
 R_{m,A}(q)=
 \bigcup_{a=1}^{A-m+1}
 \left(3^{m-1}+2^aR_{m-1,A-a}(q)\right),
 \qquad R_{1,A}(q)=\{\bar1\}\quad(A\ge1).
\]
The recurrence is for $m\ge2$, $A\ge m$; $R_{m,A}(q)=\varnothing$
when $A<m$.  For column count vectors the transition matrix is explicitly
$M_{m,a}e_r=e_{3^{m-1}+2^ar}$, with inverse
$M_{m,a}^{-1}e_s=e_{2^{-a}(s-3^{m-1})}$.  Its dependence on $a$ is
periodic modulo $\lambda=\operatorname{ord}_q(2)$, taking $\lambda=1$
for $q=1$.  Summing the periodic matrix blocks proves the rational generating
function in the companion, with every coefficient a finite sum.

For $(m,A)=(8,13)$, direct enumeration and the independent recursion agree:
\[
 R_{8,13}(233)=\mathbb Z/233\mathbb Z\setminus\{0,138\}.
\]
The $792$ packets form $99$ free cyclic classes.  Exactly $15$ minimal
representatives have numerator divisible by $7$, none has numerator divisible
by $233$, and none is integral.  The reproduced exhaustive table through
$m=16$ contains $4{,}734{,}620$ packets and zero integral packets in the eight
strip-selected labels
\[
 (5,8),(8,13),(10,16),(11,18),(13,21),(14,23),(15,24),(16,26).
\]
The resulting exclusion of nontrivial exact odd-step periods $2\le m\le16$
is valid, but it is a finite certificate rather than a novelty claim; the
published literature already contains much stronger cycle-length exclusions.

\subsection{Reproducibility and admission boundary}

The companion additionally retains the common exponent word when comparing
residues modulo coprime factors.  Its \emph{Synchronized CRT transport}
theorem gives the explicit isomorphism and transition identity, and
\emph{Information discarded by the marginals} computes the kernel of the
marginal map.  At $(m,A)=(10,16)$ the complete joint table has $3473$
occupied cells modulo $(13,499)$ but no occupied $(0,0)$ cell, although both
marginal supports are full.  All $35$ words with numerator divisible by
$499$ are squares of packets in $\Pi_{5,8}$.  Thus $499$ already rules out
an integral periodic point for every primitive word at this label.
These are independently reproduced
finite calculations, not a failure of the Chinese remainder theorem.

The root certificate
\path{certificates/affine_packet_source_and_cycle_checks.py} pins the raw
export, local antecedent, and exact primary manifestations of
B\"ohm--Sontacchi, Crandall, Lagarias, Urata, and Laarhoven--de Weger.  It
then executes the portable companion certificate
\path{research_companion/certificates/affine_packet_residue_checks.py}.
The latter uses integer and exact rational arithmetic only, checks the
packet--parity morphism on 12,089 packets, including 155,047 shortened steps,
233,525 full steps, and 78,478 odd returns.  It checks exact periods for
primitive and repeated packets, groupoid inverses and composition,
702 residue multiplicity vectors and 546 rational-series coefficient vectors,
verifies the 99-row CSV, and runs two independent full packet scans.
Within the portable \path{research_companion/} tree, the separate certificate
\path{certificates/affine_signed_necklace_checks.py} checks the two signed
fixed-content formulas on 100 parameter pairs containing
39,202 words and 39,022 primitive words, together with the cyclic freeness,
packet--binary necklace bijection, point/orbit distinction, coprime case, and
recovery of the unrestricted Moreau sum.  All three scripts reject optimized
Python.

The admitted mathematical content is therefore separated as follows.
\begin{itemize}[leftmargin=2em]
\item The rational fixed-return and parity-necklace theorems are attributed to
their published sources.
\item The packet--parity groupoid isomorphism, its fixed-content signed count,
and the displayed residue recursion are direct deductions proved in the
companion, with classical necklace lineage retained explicitly.
\item The period-five, period-eight, and through-period-sixteen statements are
finite exact calculations with pinned reproducible certificates.
\item The local technical note remains a local antecedent, not a publication
and not an authority for enlarging any published theorem.
\item No statement here proves global termination, excludes all nontrivial
cycles, or assigns a result to an author who did not state it.
\end{itemize}
